\chapter{Binomial Coefficients}

\begin{definition}
    For non-negative integers $a$, $b$, we define

    \[
    \binom{a}{b}=
    \begin{cases}
        0 \ \text{if} \ a<b\\
        \dfrac{a!}{b!(a-b)!}
    \end{cases}
    \]
    
\end{definition}

\begin{lemma}
    If $p$ is a prime, and $0<k<p$, then $p \mid \binom{p}{k}$
\end{lemma}

\section{Lucas's Theorem}

\begin{theorem}[Lucas's Theorem]
    Let $p$ be a prime and $m, n$ non-negative with base-$p$ representations $m_k p^k+m_{k-1}p^{k-1}+\dots+m_1 p+m_0=m$, and $n_k p^k+n_{k-1}p^{k-1}+\dots+n_1 p+n_0=n$ respectively. Then,

    \[
    \binom{m}{n} \equiv \binom{m_k}{n_k}\binom{m_{k-1}}{n_{k-1}}\dots\binom{m_1}{n_1}\binom{m_0}{n_0} \pmod p
    \]

\end{theorem}


